\documentclass[ecta,nameyear,draft]{econsocart}
\usepackage{amsmath,amsfonts,amssymb,amsthm,booktabs,array,longtable}
\usepackage[T1]{fontenc}
\usepackage{mathpazo,microtype,needspace,xurl}
\usepackage{xr-hyper}
\RequirePackage[colorlinks,citecolor=blue,linkcolor=blue,urlcolor=blue]{hyperref}
\startlocaldefs
\newtheorem{theorem}{Theorem}
\newtheorem{proposition}[theorem]{Proposition}
\newtheorem{lemma}[theorem]{Lemma}
\newtheorem{corollary}[theorem]{Corollary}
\newcommand{\E}{\mathbb E}
\newcommand{\R}{\mathbb R}
\endlocaldefs
\setlength{\emergencystretch}{3em}
\makeatletter\def\copyright@text{Prepared for referee review}\makeatother

\setcounter{secnumdepth}{2}

\externaldocument[][nocite]{build/supp}[supp.pdf]
\externaldocument[][nocite]{build/appendix}[appendix.pdf]
\begin{document}
\begin{frontmatter}
\title{Neural Bellman Operators}
\runtitle{Neural Bellman Operators}
\begin{aug}
\author[id=au1,addressref={add1}]{\fnms{Qian}~\snm{QI}\ead[label=e75ECTA]{qiqian@pku.edu.cn}}
\address[id=add1]{Peking University}
\end{aug}
\begin{abstract}
This paper develops Neural Bellman Operators for constructive policy evaluation and feasible improvement. Witness-preserving continuations connect approximation resources to the original economic optimum. Qualification-optimal fixed-feature readouts have exact primal--dual certificates and computable original-law moments. We derive a finite-controller displacement bound that qualifies policy-cost differences without assuming smooth acquired actors or querying a policy-value oracle. A prospective experiment evaluates four attained ReLU and Bernstein controllers under the original nonlinear investment law, with every path independently reconstructed by integer interval arithmetic. Both ReLU controllers have lower expected cost than their resolution-matched Bernstein comparators. Seven of eight paired intervals meet the predeclared precision target; the remaining failure is reported. The identified gains yield an explicit priced-adoption region. Predictable qualification, paid witness reuse, complete resource accounts and prior adverse comparisons distinguish policy quality, verification accuracy and computational work.
\end{abstract}
\begin{keyword}
\kwd{Dynamic programming}\kwd{Neural approximation}\kwd{Policy certification}
\kwd{Information acquisition}\kwd{Computational economics}
\end{keyword}
\end{frontmatter}
\section{Introduction}\label{sec:introduction65}
A numerical continuation matters because it helps produce an economic decision. The error of a fitted value, the loss of the implemented policy and the resources required to verify that loss are different objects. A policy may improve without a tighter numerical certificate; a tighter certificate may leave the policy unchanged. A useful computational method must connect its approximation and resource choices to a feasible decision under the original economic law.

This paper develops Neural Bellman Operators (NBOs) along that chain. The constructive implementation evaluates a continuation against its own fitted future and retains feasible action witnesses. Its exact affine--ReLU realization and primitive error budget do not require a theorem that unrestricted nonconvex training succeeds. A complementary recursive implementation queries the original Bellman problem only where action search requires information. Trained readouts can propose actions, but verified lower and upper bounds decide whether the resulting policy meets its original-optimum contract. These are distinct uses of a learned function, with different costs and comparison burdens.

The investment economy has state-dependent capacity, bilinear transitions and continuous innovations. Primitive curvature controls the convexity defect of the transition and gives regularity of the true Bellman value. A strong-convexity upper chord can certify a rounded proposal before another endpoint query; protected conventional refinement remains available otherwise. This common certificate applies to neural and conventional proposals. Its validity is not itself a neural advantage. The operational and two-control results are retained in Sections~\ref{sec:operational67} and \ref{sec:vector67} of the Numerical and Development Appendix.

The central additional role of the neural object is as an attained continuation that determines the complete controller. A frozen fitted continuation supplies an exact original-law evaluation identity with a concentration range governed by verified Bellman-residual oscillations and implemented greedy gaps. Its terminal correction is retained when the fitted terminal value is inexact, and state-independent shifts cancel. For shallow affine--ReLU continuations, one- and two-innovation conditional expectations are finite positive-part moment calculations. Explicit state-and-action covers qualify the residual without evaluating the optimal Bellman descendant tree. The uniform cover and the fit remain paid operations.

The attained fitting criterion is chosen for that qualification problem. A finite linear program minimizes the verified residual-oscillation bound over a fixed feature class, conditional on the service's actual fitted future. Exact rational primal--dual calculations certify the achieved objective gap. Backward construction gives a fitted-continuation-to-policy-bound chain without substituting a solver message for a proof. Matched ReLU and Bernstein constructions have nine coefficients, the same coefficient limits and the same economic information and cover precision. The result concerns those attained finite readout problems, not global optimization of hidden-layer parameters or joint multi-date minimization of true regret.

A uniform certificate still does not rank actual costs. We complete the economic comparison by directly evaluating the four returned controllers under the original continuous law. The sample and all estimands are fixed after the controller identities are frozen and before observing their costs. A new displacement theorem exploits the full finite action tables: cell pairs close enough to contain coupled states bound the next action difference, and primitive transition and payoff moduli propagate that information into a path-cost support. The actors may be discontinuous. The support is computed before sampling, includes boundary cases and does not use a policy-value oracle.

This produces an informative comparison in the original two-date nonlinear economy. Both attained ReLU controllers have lower expected cost than their resolution-matched Bernstein controllers. Refinement lowers expected cost within each basis. The two refinement differences are also equivalent at the separately declared positive tolerance. Seven of the eight paired intervals meet the predeclared total-width target; the remaining precision failure has a strict sign and is reported without additional sampling. Each of the 655,360 controller-path evaluations is independently reconstructed by arbitrary-precision interval arithmetic over its entire continuous bin. The thirteen-estimand family uses 131,072 shared rows, not five independent datasets.

The distinction between the qualification objective and expected cost is observable. At the coarser resolution the Bernstein continuation has a smaller uniform qualification bound, but its controller has higher expected cost. The former is a bound against the original optimum over the entire state domain; the latter is an average under the original initial law. A criterion can be useful for certifying a policy without supplying its economic ranking. The new signed gains therefore extend, rather than restate, the attained fitting theorem.

The policy-cost term makes a purchase decision possible. The gain interval is combined with the value of a unit of normalized loss, the priced one-off construction and qualification work, the installation charge and the intended number of deployments. A positive lower net benefit supports adoption; a negative upper net benefit rules it out. The same simultaneous cost event supports the full parameter region. Processor seconds are not treated as monetary units, and the inference sample count is not the number of future deployments over which setup can be amortized. The model's prices remain explicit external inputs rather than an invented empirical calibration.

Two other resource problems retain separate solutions. Predictable continuation and support choices yield a cost certificate valid under optional stopping, with fixed-controller costs distinguished from conditional averages of changing controllers. A lower original-Bellman bound and its feasible upper witness can also be transported to a nearby state after paying for state displacement, capacity repair and action rounding. A failed check invokes the original verifier. The finite-cover theorem specifies the tighter fill precision and retention needed to bound expensive fills. Neither a cache hit nor the absence of a state lattice removes the cost of a first descendant tree.

The earlier scalar, vector, validation and sixty-four-block reuse studies remain part of the paper. They retain both favorable small-task reuse observations and adverse cold-start and difficult-task comparisons. Their twelve centered cost intervals establish a declared equivalence margin, not a strict ranking. The new table-controller ranking does not reverse those different service observations or establish superiority to an unevaluated adaptive controller. Likewise, the two-date comparison does not establish a long-horizon or high-dimensional speed result. These distinctions identify the contribution of the attained neural continuation without attributing every successful protected return to neural fitting.

\paragraph{Relation to existing methods.}
Lipschitz extension and distance transforms underlie the constructive envelope \citep{mcshane1934,felzenszwalb2012}. Approximate and safe policy iteration have established error and performance analyses \citep{munos2003,pirotta2013,scherrer2015}; \citet{yang2020convex} studies convex dynamic-programming approximation. Learned warm starts and optimization-based control have their own literature \citep{sakaue2023warm,sambharya2023,agrawal2020control}. Our contribution connects the constructed continuation and these numerical choices to an acquired feasible policy, a verified stopping decision and the complete cost of returning and qualifying that policy.

Value-function control variates and telescoping identities are established evaluation tools \citep{munos2006variance,jiang2016}. Here a verified residual-oscillation bound supplies both an original-optimum policy account and an original-law concentration range. Exact conditional ReLU integration and uniform qualification specify what is purchased when optimal-value centering is replaced by a fitted continuation. The finite-controller displacement result supplies a complementary range directly from discontinuous action tables. The endpoint empirical-Bernstein calculation uses the bounded-sample inequality of \citet{maurer2009}; we derive the numerical support and identify its actual policy estimand, rather than claim a new concentration inequality.

Time-uniform concentration through nonnegative supermartingales is established \citep{howard2020time}; off-policy confidence sequences have been developed for contextual-bandit evaluation \citep{karampatziakis2021}. The predictable cost theorem uses those principles for original-law NBO contrasts enclosed over continuous innovation bins. Its target and reset contract differ from an adaptive average of changing policies. Fixed-sample comparisons below do not acquire optional-stopping validity merely because a sequential theorem appears elsewhere in the paper.

\paragraph{Reading path.}
Sections~\ref{sec:core49}--\ref{sec:query63} give the original constructive operator, investment primitives and recursive guarantee. Sections~\ref{sec:frozen72}--\ref{sec:sequential73} develop fitted-continuation evaluation, attained readout qualification and predictable cost certification. Section~\ref{sec:displacement75} gives the new finite-controller support and adoption account. Paid witness transport and the preceding construction and reuse studies retain their own roles. Section~\ref{sec:study75} reports every new original-law comparison. Complete primitive and inferential details are in the Technical Supplement; the operational, vector, centering and prior economic sections remain compiled in the Numerical and Development Appendix. The broader diffusion, preference, temporal-self and game applications remain in their unchanged companions under their original hypotheses.

\section{The economic operator and its policy account}\label{sec:core49}
Consider dates $t=0,\ldots,T$ and an invariant compact state domain $X\subset\R^d$. At date $t$, the action belongs to a nonempty compact set $A_t(x)$. The transition is $F_t(x,a,Z_{t+1})$, the current cost is $c_t(x,a)$, and the terminal cost is $g(x)$. Write $B_0=1$ and $B_t=\prod_{j<t}\beta_j$, with $\beta_j>0$. An implemented policy includes its observation rule, action selection and feasible numerical realization, not just its stored parameters.

For a continuation $v$, define
\begin{equation}\label{eq:operator49}
 Q_t(v;x,a)=c_t(x,a)+\beta_t\rho_t(v(F_t(x,a,Z))),\qquad
 \mathcal T_tv(x)=\inf_{a\in A_t(x)}Q_t(v;x,a).
\end{equation}
Throughout the deterministic comparison, the evaluations $\rho_t$ are monotone and cash invariant on a common domain containing all functions and constant translates used below. They are therefore nonexpansive in the supremum norm. Continuous objectives, effective compact feasible correspondences and measurable minimizers, or measurable arbitrarily near-minimizers with their error charged, are assumed. The stipulated conditional recursion gives Bellman comparison against all admissible adapted controls. In the experiments $\rho_t$ is conditional expectation, costs are bounded and innovations have their original continuous uniform law. The direct-cost inference in Section~\ref{sec:decisions49} is specific to expectation; it is not automatically an inference theorem for nonlinear recursive utility or games.

A Neural Bellman Operator is a construction and evaluation procedure for \eqref{eq:operator49}. It forms targets using its \emph{own} fitted $f_{t+1}$, returns a neural continuation and a feasible policy, and records an economic error allowance. It does not receive the optimal continuation as a training label. Arbitrary trained networks can be evaluated by the first result below. The subsequent constructive backend specifies how a network and its actor are obtained without assuming successful nonconvex optimization.

\subsection{Action-dependent errors and complete-policy loss}
Let $H(a)$ be a true current cost and $\widehat H(a)$ its continuation approximation. If $\widehat a$ is $\delta$-optimal for $\widehat H$ on the same feasible action set, then
\begin{equation}\label{eq:center49}
 H(\widehat a)-\inf_aH(a)\le\delta+
 \operatorname{osc}_{a}\{\widehat H(a)-H(a)\}.
\end{equation}
Indeed, compare $\widehat a$ with an arbitrarily near-optimal true action and cancel the approximate objectives. The infimum limit gives the result. Additive continuation levels cancel; action-dependent discrepancies do not. Under a change in future policies, the continuation must be transported to the new economic problem before this inequality licenses reuse. The full development edition retains the original quantitative transport, recursive-domain and economic-application results. No transport bound is presumed valid after an unaccounted change of constraints or preferences.

\begin{proposition}[One-sided policy sandwich]\label{prop:sandwich49}
Suppose, uniformly in state,
\[
 f_t-\mathcal T_tf_{t+1}\le u_t,\qquad
 \mathcal T_t^\pi f_{t+1}-f_t\le v_t,
 \qquad \ell_T\le f_T-g\le u_T.
\]
Then the actual feasible policy satisfies
\begin{equation}\label{eq:sandwich49}
 0\le J_0^\pi-V_0^*\le
 \sum_{t<T}B_t(u_t+v_t)+B_T(u_T-\ell_T).
\end{equation}
The bound is invariant to date-specific constant shifts of the fitted continuations, provided the associated residual accounts are shifted consistently.
\end{proposition}
\noindent The proof is in Supplement, Section~\ref{supp:proof-core49-1}.


A sharper valid residual account for the same continuation and the same deployed actor sharpens knowledge of the policy. It does not retrain that policy or change its realized cost. Conversely, subtracting two upper bounds in \eqref{eq:sandwich49} does not identify the difference in actual costs. These distinctions govern both the historical comparisons and the new experiments.

\section{Constructing and compiling a feasible neural policy}\label{sec:construction49}
Assume the primitive objective has, for feasible state--action pairs, state modulus $A_t^Q=C_{x,t}+\beta_tF_{x,t}L_{t+1}$ and action modulus $D_t=C_{a,t}+\beta_tF_{a,t}L_{t+1}$. Let
\[
 d_H(A_t(x),A_t(z))\le K_t^A\|x-z\|_1,
 \qquad L_t=A_t^Q+D_tK_t^A,\qquad L_T=L_g.
\]
The Hausdorff condition and the primitive moduli imply that $\mathcal T_tf_{t+1}$ is $L_t$-Lipschitz whenever $f_{t+1}$ is $L_{t+1}$-Lipschitz. All moduli and numerical integration procedures are supplied by the economic primitives, not fitted from successful outcomes.

At state nodes $x_{t,i}$ with covering radius $h_t$, form effective feasible $k_t$-nets. Evaluate their own-future objectives with certified error at most $e_t$. Minimize the computed labels and retain the corresponding \emph{original} action $a_{t,i}$. A Borel repair $\Gamma_{t,i}(x)$ is feasible at $x$ and satisfies
\begin{equation}\label{eq:repair49}
 \|\Gamma_{t,i}(x)-a_{t,i}\|_1\le K_t^A\|x-x_{t,i}\|_1+\zeta_t.
\end{equation}
The allowance $\zeta_t$ concerns displacement; feasibility itself may not fail. Define
\begin{equation}\label{eq:cone49}
 f_t(x)=\min_i\{y_{t,i}+L_t\|x-x_{t,i}\|_1\},\quad
 i_t(x)=\min\arg\min_i\{y_{t,i}+L_t\|x-x_{t,i}\|_1\}.
\end{equation}
Return $\pi_t(x)=\Gamma_{t,i_t(x)}(x)$. Terminal labels approximate $g$ with error $e_T$ and use the same cone construction with modulus $L_g$.

\begin{theorem}[Own-future construction with feasible witnesses]\label{thm:construct49}
The functions in \eqref{eq:cone49} have exact affine--ReLU realizations and Lipschitz moduli $L_t$ independent of label error. The returned policy is feasible and measurable. It satisfies
\begin{align}\label{eq:enclosures49}
 -e_t\le f_t-\mathcal T_tf_{t+1}&\le D_tk_t+e_t+2L_th_t,\\
 \mathcal T_t^\pi f_{t+1}-f_t&\le e_t+D_t\zeta_t,\notag
\end{align}
and hence
\begin{equation}\label{eq:constructionloss49}
 J_0^\pi-V_0^*\le
 \sum_{t<T}B_t(D_tk_t+2e_t+2L_th_t+D_t\zeta_t)
 +B_T(2e_T+2L_gh_T).
\end{equation}
With effective covers, repairs and certified queries at the prescribed accuracy, allocating each nonnegative term a share of a positive economic tolerance gives a terminating finite construction. A fixed numerical precision or a capped finite experimental ladder need not contain that allocation.
\end{theorem}
\noindent The proof is in Supplement, Section~\ref{supp:proof-core49-2}.


The Lipschitz-extension principle is classical \citep{mcshane1934}. A native min-plus envelope and its exact ReLU circuit, using the same labels, witnesses and ties, are \emph{one method with two encodings}. They return the same function and policy. Neither Theorem~\ref{thm:construct49} nor a faster representation proves a uniquely neural capability or a training-success probability for Adam or L-BFGS. The mathematical increment is the complete feasible, own-future construction and implementation account; the empirical question is its work and policy value relative to conventional generators.

\subsection{Original-owner compilation on a tensor cover}
On a tensor cover $\mathcal X=\prod_{j=1}^d\{x_{j,0},\ldots,x_{j,n_j}\}$, put $S=\prod_j(n_j+1)$. For arbitrary labels $y_i$, including inconsistent inexact labels, compute at each grid node $v$
\[
 (z_v,o_v)=\min_{\rm lex,i}(y_i+L\|v-x_i\|_1,i).
\]
The index $o_v$ is the original label owner, not the transformed node $v$.

\begin{theorem}[Continuous original-owner identity]\label{thm:compile49}
For a closed tensor cell $C$ with corners $\mathcal V(C)$, and $x\in C$,
\begin{align}\label{eq:compile49}
 f(x)&=\min_{v\in\mathcal V(C)}(z_v+L\|x-v\|_1),\\
 (f(x),i_*(x))&=\min_{\rm lex,v\in\mathcal V(C)}
 (y_{o_v}+L\|x-x_{o_v}\|_1,o_v).\notag
\end{align}
Coordinatewise forward and backward relaxations compute $(z,o)$ in $O(dS)$ additions and comparisons and $O(S)$ storage. After locating the cell, a query uses $O(d2^d)$ arithmetic operations. Exact rational comparisons preserve all ties, including $L=0$. Using the original owner's repaired action preserves the uncompiled policy, its cost and its certificate exactly.
\end{theorem}
\noindent The proof is in Supplement, Section~\ref{supp:proof-core49-3}.


The nodal algorithm is the classical distance transform \citep{felzenszwalb2012}; preserving the original action through continuous off-grid queries is the policy obligation. If there are $J_{t,i}$ node actions and $M_t$ integration representatives, write $Q_t=M_t\sum_iJ_{t,i}$. Literal dense cone evaluation costs $O(\sum_t dS_{t+1}Q_t)$ in representation work. Compilation replaces this by $O(\sum_t[dS_{t+1}+Q_td2^d])$, plus cell location, primitive evaluation, integration, feasible repair, certificate and output. Rational bit lengths and state boxes crossing many cells incur additional costs. There is no dimension-free complexity assertion.

\subsection{Acquired states and conventional actors}
A $b$-bit coordinate sensor reports a dyadic cell and its midpoint $\widehat x$. On the unit cube, $\|x-\widehat x\|_1\le r(b)=d2^{-b-1}$. For the investment capacity $\kappa(x)$ below, clip the original action to the smallest capacity throughout that cell and round down to an action quantum $\Delta_t$. This is feasible at every state in the cell. Relative to an exact true-state repair, its displacement is at most $2K_t^Ar(b)+\Delta_t$. Transporting the witness cone from $\widehat x$ to $x$ adds at most $2L_tr(b)$, plus any certified selector objective error $\nu_t$. Thus the implemented-policy account adds
\begin{equation}\label{eq:acquire49}
 \sum_{t<T}B_t\{(2L_t+2D_tK_t^A)r(b)+D_t\Delta_t+\nu_t\}.
\end{equation}
No continuity of the discrete owner selector is assumed. At uncertain numerical comparisons the simulator retains all possible actors, rather than choosing a favorable one.

For conventional FVI the continuation is multilinear on its realized tensor grid. Its actual modulus is bounded by the largest coordinate edge slope. A nearest-node actor uses a feasible node minimizer and true-state repair. Let $h_t$ be the largest half-cell $\ell^1$ radius, $L_t^{\rm gr}$ the Bellman graph modulus, and $e_t$ the node query allowance. Its optimal residual upper bound is $D_tk_t+e_t+L_t^{\rm gr}h_t$. Its selected-policy upper bound is $e_t+E_t$, where
\[
 E_t=\sup_x\{y_{i(x)}+L_t^{\rm gr}\|x-x_{i(x)}\|_1-f_t(x)\}.
\]
Partitioning cells by nearest-node boundaries makes the expression separately affine in each coordinate. Its maximum is therefore at a partition vertex; the implementation encloses all such vertices. Both uniform and curvature-refined FVI receive this sharper one-sided account, not an artificially doubled baseline. Their acquired-state allowance is derived separately in \eqref{eq:fvisensor49}.

\section{Original-optimum accuracy in the investment economy}\label{sec:accuracy62}
The exact witness theorem answers an implementation question conditional on a continuation. An economist also needs to know how close the complete implemented policy is to the original optimum. We now establish that connection in the investment economy itself. The argument exploits a property not used by the earlier cell-range calculation: despite the bilinear transition, the Bellman objective remains jointly convex. This permits a second-order interpolation account and transfers verified nodal decisions to an implementable policy over the entire state domain. It does not assume that a fitted neural continuation is convex.

\paragraph{Relation to the preceding two-control query.}
The earlier query experiment used the fixed shared capacity $1/8$. The complete dynamic extension below uses the original state-dependent capacity $c(x)$, with the retained exposure columns, action-cost interaction and second shock. Its nesting claim concerns the original scalar economic model; it is not an assertion that a full dynamic policy was already present in the earlier isolated query.

\subsection{Primitives and Bellman regularity}
Let $d\ge2$ be even, let indices be cyclic, and set $\beta=15/16$. Write
\begin{align*}
 S(x)&=\frac2d\sum_j(x_j-5/8)^2+
       \frac1{4d}\sum_j(x_j-x_{j+1})^2+
       2\left(\frac12-\frac2d\sum_jx_j\right)_+^2,\\
 g(x)&=S(x)+\frac2d\sum_j(x_j-5/8)^2,\qquad
 c(x)=\frac18+\frac1{8d}\sum_jx_j.
\end{align*}
The original scalar control has $0\le a\le c(x)$, running cost $S(x)+a^2+4a^4$, and transition
\[
 F_j(x,a,z)=\frac1{16}+\frac{x_j}{2}+\frac{x_{j+1}}8+
       \frac{x_j(1-x_{j+1})}{16}+B_ja+s_jz,
\]
where $(B_j,s_j)=(1/2,1)$ at alternating coordinates and $(1/4,-1)$ at the others, and $z$ is uniform on $[-1/32,1/32]$.

For the retained two-control extension, $a_1,a_2\ge0$, $a_1+a_2\le c(x)$, the second exposure column interchanges $1/2$ and $1/4$, and the action cost is
\[
 R(a)=a_1^2+a_2^2+4(a_1^4+a_2^4)+a_1a_2/4.
\]
An independent common-direction uniform innovation of radius $1/64$ is added. This is the previously specified extension, nested at $a_2=0$ when that innovation is disabled. Both models preserve $[0,1]^d$. Let $m\in\{1,2\}$ denote the number of controls and define
\begin{align*}
 Q_t^*(x,a)&=S(x)+R(a)+\beta\mathbb E V_{t+1}^*(F(x,a,Z)),\\
 V_t^*(x)&=\min_{a\in A(x)}Q_t^*(x,a),\qquad V_T^*=g.
\end{align*}
All costs and laws in this display are the economic primitives, not a fitted surrogate.

\begin{theorem}[Regularity of the original Bellman problem]\label{thm:regularity62}
For either control specification and every finite horizon, $V_t^*$ is convex and continuously differentiable in the interior, with the corresponding one-sided extension on the cube. It is $G_t$-Lipschitz in the $\ell^1$ norm and $M_t$-semiconcave in the Euclidean norm, where valid constants are
\begin{align*}
 G_T&=10/d,&G_t&=\frac{1037}{128d}+\beta\frac{47}{64}G_{t+1},\\
 M_T&=26/d,&M_t&=\frac{22}{d}+\frac{21}{256d}
             +\beta\left(\frac9{16}M_{t+1}+\frac18G_{t+1}\right).
\end{align*}
In particular, $G_t\le27/d$ and $M_t\le54/d$. The function $Q_t^*$ is jointly strongly convex on its convex feasible graph, with state curvature at least $107/(128d)$ and action curvature at least $7/4$.

A valid semiconcavity constant along each coordinate line is
\begin{align*}
 D_T&=\frac9d+\frac{16}{d^2},\\
 D_t&=\frac5d+\frac{16}{d^2}+\frac{21}{256d^2}
 +\beta M_{t+1}\left(\frac{85}{256}+\frac{23}{256d}\right).
\end{align*}
Consequently, for the positive multilinear interpolant $I_n$ on a mesh of spacing $1/n$,
\begin{equation}\label{eq:interpolation62}
 0\le I_nV_t^*(x)-V_t^*(x)\le\kappa_t,
 \qquad \kappa_t=\frac{dD_t}{8n^2}.
\end{equation}
\end{theorem}

\noindent The proof is in Supplement, Section~\ref{supp:proof-accuracy62-print-1}.


The theorem is a structural statement about these economic primitives. It is not a claim that every controlled diffusion, recursive preference specification, or game in the complete development edition has a convex Bellman problem. Nor does it require a signed-weight neural fit to preserve convexity. Convex approximation of dynamic programs has an established literature; the additional work here is to verify the needed curvature for the original nonlinear capacity-constrained model and connect it to the actual NBO actor and complete error account.

\subsection{Certified nodal Bellman bounds}
Suppose nodal arrays satisfy $l_{t+1}\le V_{t+1}^*\le u_{t+1}$ at all grid vertices. Equation~\eqref{eq:interpolation62} gives the functions $I_nl_{t+1}-\kappa_{t+1}$ and $I_nu_{t+1}$ as valid lower and upper continuations. The lower function need not itself be convex; convexity is used for the unknown true value and for the comparison argument, not asserted of arbitrary endpoint arrays.

Divide each uniform innovation component into $q$ equal bins and use its conditional mean. Convexity gives a lower quadrature bound for the true continuation. Semiconcavity gives an upper error
\begin{equation}\label{eq:quadrature62}
 \zeta_{t+1}=\frac{M_{t+1}}{6q^2}
                      \sum_{r=1}^{s}\rho_r^2\|s_r\|_2^2,
\end{equation}
where $s=1$ or $2$, $\rho_1=1/32$, and $\rho_2=1/64$ in the extension. The expectation remains the original continuous law. At the terminal date $g$ is evaluated directly, so no terminal interpolation allowance is needed.

At a vertex $v$, let $\mathcal A_A(v)$ be the capacity simplex lattice with spacing $c(v)/A$, including all its faces. For $m=1,2$, a minimizer can be represented by neighboring lattice vertices whose coordinate variances sum to at most $m c(v)^2/(4A^2)$. A valid action smoothness constant is
\[
 H_t=\frac{21}{4}+\beta M_{t+1}\|B\|_2^2,
 \quad \|B\|_2^2=\frac{5d}{32}\ (m=1),\quad
 \|B\|_2^2=\frac{9d}{32}\ (m=2).
\]
Taylor's inequality and that barycentric representation imply
\begin{equation}\label{eq:action-cover62}
 0\le\min_{a\in\mathcal A_A(v)}Q_t^*(v,a)-V_t^*(v)
       \le\delta_t:=\frac{mH_t}{128A^2}.
\end{equation}
The linear Taylor terms cancel even when the optimum lies on a face. This is an error against the continuous feasible optimum, not only against the finite action menu.

Let $q_t^-(v,a)$ and $q_t^+(v,a)$ include all interpolation, quadrature and arithmetic errors and enclose $Q_t^*(v,a)$. Then
\begin{equation}\label{eq:nodal62}
 l_t(v)=\max\{0,\min_{a\in\mathcal A_A(v)}q_t^-(v,a)-\delta_t\},
 \qquad
 u_t(v)=\min_{a\in\mathcal A_A(v)}q_t^+(v,a)
\end{equation}
are valid nodal Bellman bounds. The zero truncation uses the nonnegative economic costs. Extra neural or conventional proposals may improve the upper side but do not remove the lower action-cover allowance. Positive interpolation, averaging and minimum maps are monotone and nonexpansive; the separate allowances in \eqref{eq:interpolation62}--\eqref{eq:action-cover62} are therefore propagated, rather than being discarded because an endpoint calculation contains the truth.

\subsection{From a fitted witness to the complete implemented policy}
Let $a_t(v)\in A(v)$ be the actual action returned at each vertex by any generator. It may be selected by the exact neural witness, a conventional fitted continuation, or their common verified fallback. Write $w_v(x)$ for the multilinear weights and set $\bar\pi_t(x)=\sum_vw_v(x)a_t(v)$. Affineness of capacity makes this action feasible. Let $q_t^{\mathrm{sel},+}(v)$ be a verified upper bound on its original-optimum continuation value $Q_t^*(v,a_t(v))$, using the same lower and upper Bellman reference for all generators.

\begin{theorem}[Actual-policy, whole-domain certificate]\label{thm:policy62}
Set
\[
 e_t=\max_v\{q_t^{\mathrm{sel},+}(v)-l_t(v)\}+\kappa_t.
\]
Suppose the implemented policy $\widehat\pi_t$ is feasible and differs from $\bar\pi_t$ by at most $b_t$ in action $\ell^1$ norm, uniformly over the state cube. Define
\[
 \Lambda_t=\frac{13}{16}+\beta\frac{3d}{8}G_{t+1},\qquad
 D_T^\pi=0,\quad D_t^\pi=e_t+\Lambda_tb_t+\beta D_{t+1}^\pi.
\]
Then, at every true state and date,
\[
 0\le J_t^{\widehat\pi}(x)-V_t^*(x)\le D_t^\pi.
\]
This statement is conditional on the actual recorded policy, not on success of its fitting algorithm. It applies equally to a pure fitted policy and to a policy augmented by verified common candidates.
\end{theorem}

\noindent The proof is in Supplement, Section~\ref{supp:proof-accuracy62-print-2}.


In the execution, observation is rounded down at forty bits. An acquired state $y\le x$ has no greater capacity than the true state. Each interpolated action is then rounded down at twenty bits using the lower endpoint of its arithmetic enclosure. If $L_{ij}$ bounds the derivative of action coordinate $i$ with respect to state coordinate $j$, and $\eta$ bounds native interpolation error, a valid implementation allowance is
\[
 b_t\le 2^{-40}\sum_{i,j}L_{ij}
           +m(2^{-20}+2\eta+\eta_{\mathrm{out}}).
\]
Here $L_{ij}$ is computed from the maximum adjacent nodal difference times $n$, and $\eta_{\mathrm{out}}$ covers endpoint subtraction. All these numbers are recorded outward. This new barycentric acquired actor is explicitly distinguished from the earlier discontinuous cell-constant actor; an old actor is not silently replaced during a purported fixed-policy revalidation.

\subsection{Arithmetic, work, and scope}
The native interpolation kernel forms nonnegative tensor weights and a weighted sum of at most $2^d$ stored binary64 endpoints. It is compiled without reassociation or floating-point contraction. For $K=(3d+4)2^d+32$, an explicit conservative absolute allowance is
\[
 \eta=\frac{K2^{-53}}{1-K2^{-53}}\max_v|z_v|+K2^{-1022}.
\]
The final term conservatively covers underflow. Input coordinates, affine indexing, finite endpoints and the binary64 format are checked. Roundoff expansion and the ordinary interval cost arithmetic remain separate from the analytic approximation bounds. Exact rational fixtures compare the native sum with its mathematical tensor interpolant.

With $N=(n+1)^d$ vertices, $K_A=A+1$ scalar actions or $K_A=(A+1)(A+2)/2$ two-control actions, the reference calculation uses $O(TNK_Aq^s2^d d)$ elementary work under the explicit tensor implementation. Stored reference values use $O(TN)$ words and a policy uses $O(TNm)$ words. Streaming batches avoid a state-pair table. Proposal fitting, exact rational witness work, actor transfer, serialization, independent cost inference and compilation are additional costs, not absorbed into this bound. The analytic errors are second order in state, action and innovation resolution; the state-count and tensor-evaluation factors still depend exponentially on dimension. The result supplies a matched original-optimum comparison without asserting that these costs have disappeared.

\section{Neural Bellman queries without a state lattice}\label{sec:query63}
The preceding construction makes numerical error accountable, but a complete tensor reference can absorb the benefit of a fitted continuation. We therefore change the computational organization, not the economic problem. A Neural Bellman selector proposes an action at the state actually acquired. A recursive interval query then bounds the original Bellman value at that state and checks the proposed action against the entire continuous feasible interval. No conventional policy is precomputed at all states. The controller is the resulting callable algorithm, not the unverified proposal alone.

The distinction matters for attribution. Learning may reduce the work needed to locate a useful action. It does not supply the inequalities used to certify that action. The adaptive conventional comparator receives the same inequalities and the same original-optimum target. The trained value-fit policies in the preceding study remain distinct from the action-query predictor introduced here.

\subsection{A curvature lower bound on an adaptive action partition}
Use the original scalar economy of Section~\ref{sec:accuracy62}, with $r$ dates remaining. Write $V_r$ for its optimal value and $Q_r(x,a)$ for the value of a first action followed by the optimal future. Let $G_r,M_r$ be the primitive constants in Theorem~\ref{thm:regularity62}, indexed by remaining horizon, and put
\[
 H_r=\frac{21}{4}+\beta M_{r-1}\frac{5d}{32},
 \qquad \Lambda_r=\frac{13}{16}+\beta G_{r-1}\frac{3d}{8}.
\]
Thus $H_r$ is an upper curvature bound in the scalar action, and $\Lambda_r$ bounds the absolute action derivative. Neither depends on a fitted model.

\begin{lemma}[Parabolic minorants]\label{lem:parabolic63}
Suppose $f$ is $H$-semiconcave on $[l,u]$, and $L_l\le f(l)$, $L_u\le f(u)$. For $s\in[0,1]$,
\[
 f(l+s(u-l))\ge (1-s)L_l+sL_u
                 -\frac H2(u-l)^2s(1-s)=:p(s).
\]
Consequently $\min_{[0,1]}p$ is a lower bound for the minimum of $f$ on the interval. With $K=H(u-l)^2/2$ and $D=L_u-L_l$, this lower bound is $L_l$ if $D\ge K$, $L_u$ if $D\le-K$, and otherwise
\[
 \frac{L_l+L_u}{2}-\frac K4-\frac{D^2}{4K}.
\]
The endpoints $L_l,L_u$ need not themselves belong to a convex interpolant. Increasing $K$ preserves the lower-bound property.
\end{lemma}
\noindent The proof is in Supplement, Section~\ref{supp:proof-query63-1}.


At an acquired state, retain an ordered partition of the feasible action interval and an outward value bracket at every partition point. The minimum of the parabolic minorants bounds the continuous-action infimum from below. The smallest upper endpoint supplies an actual feasible witness. A neural or polynomial proposal inserts another point into this partition; it cannot remove the endpoints or any part of the feasible interval. The running lower bound is the maximum of all previously valid global lower bounds, so new information cannot weaken the certificate. The upper bound is the minimum over every evaluated feasible witness.

\subsection{Recursive queries under the original continuous law}
For an endpoint action with $r>1$, divide the scalar innovation into $q_r$ equiprobable bins and query $V_{r-1}$ at each conditional-mean next state. If each child returns width at most $\eta_r$, the endpoint bracket has width at most
\begin{equation}\label{eq:querywidth63}
 \beta\eta_r+\frac{\beta M_{r-1}d}{6144q_r^2}+\omega_r.
\end{equation}
Here $\omega_r$ includes outward primitive arithmetic and the Lipschitz expansion from an interval-valued next state to its queried center. Convexity provides the lower conditional-mean inequality; semiconcavity and the exact conditional variance provide the upper allowance. The innovation remains continuous. There is no state-interpolation term because a child value is queried rather than read from a state lattice.

For a requested local gap $\delta$, the implementation sets $\eta_r\le\delta/(4\beta)$ and chooses a dyadic $q_r$ so that the second term in \eqref{eq:querywidth63} is at most $\delta/4$. A separate endpoint-width guard charges arithmetic and the small capacity-rounding remainder. Failure of that guard is a precision failure, not a successful return. The terminal query uses analytic integration instead of recursive quadrature.

\begin{theorem}[A state-grid-free NBO service]\label{thm:query63}
Consider the original scalar investment economy. At every recursive query, suppose the outward endpoint and capacity-transfer accounts give width at most $5\delta/8$. Let the lower bound be the running maximum described above. Refine an interval attaining the smallest current minorant whenever the global upper-minus-lower gap exceeds $\delta$. Each new point must lie between one fifth and four fifths of that interval. Arbitrary additional learned proposals may be inserted initially. Then:

(i) every returned interval $[L_r(x),U_r(x)]$ contains $V_r(x)$, and its returned feasible action $a_r(x)$ satisfies
\[
 Q_r(x,a_r(x))-V_r(x)\le U_r(x)-L_r(x)\le\delta;
\]
(ii) the real-arithmetic refinement terminates. With at most one interior learned proposal, a sufficient point budget is
\begin{equation}\label{eq:querycap63}
 6+\left\lceil5c(x)\sqrt{H_r/(3\delta)}\right\rceil;
\end{equation}
(iii) for the policy that calls this service at every subsequent date, the original-law loss satisfies
\begin{equation}\label{eq:querypolicy63}
 J_T^{\pi}(x)-V_T(x)
 \le\sum_{t=0}^{T-1}\beta^t
 \left[\delta_t+
 \left(8+\frac{11\beta dG_{T-t-1}}{16}
                  +dG_{T-t}\right)2^{-b}\right],
\end{equation}
when the state is acquired downward at $b$ bits per coordinate and every returned action is feasible at that acquired state. All action rounding must precede its endpoint evaluation. The inequalities hold uniformly in the true initial state, not only at the states used in an experiment.
\end{theorem}
\noindent The proof is in Supplement, Section~\ref{supp:proof-query63-2}.


The finite-precision contract is explicit. The executed service uses the width guard, a 512-point cap, and a 32-bit downward action lattice. A precision or cap failure withholds the floating certificate and invokes a separately counted exact-rational enumeration. This fallback constructs values only at the required descendants; it does not load a precomputed conventional policy. The rational fallback chooses dyadic action and quadrature counts until each cover or quadrature allowance is at most one quarter of the request, and requests child width at most one quarter divided by the discount. Provided $\delta>8\Lambda_r2^{-32}$ at every recursive level, its exact brackets have width at most three quarters of the request. This proves a finite return for every acquired state in the declared tasks, including states absent from the workload. Rational bit-operation cost is additional work and is not replaced by a unit-cost floating clock. The mathematical point budget is checked against the declared tolerance range; termination in an arbitrary floating implementation is not inferred merely from the real-arithmetic argument. Fitting convergence is unnecessary for validity: any finite prediction, and even its clipped replacement when numerically invalid, is only a proposed query.

\subsection{A certified analytic terminal kernel}
The last decision admits an especially strong common comparator. Let $y=F(x,0,0)$, $b_0=1/2-2\operatorname{mean}(y)$, and $\gamma_j=B_j$. Define
\[
 D(x)=\frac8d\sum_j\gamma_j(y_j-5/8)
   +\frac1{2d}\sum_j(\gamma_j-\gamma_{j+1})(y_j-y_{j+1}).
\]
For every even $d$, the exact terminal action derivative is
\begin{equation}\label{eq:cubic63}
 16a^3+\left(2+\frac{41\beta}{32}\right)a
 +\beta D(x)-3\beta(b_0-3a/4)_+.
\end{equation}
Each branch is a strictly increasing cubic. Its floating closed-form root supplies only a candidate. The independently evaluated outward derivative and exact-law objective certify that candidate; a safeguarded derivative bracket remains available if the candidate is inaccurate.

For completeness, if $f$ is $\mu$-strongly convex and an outward interval encloses $f'(a)$, the largest compatible projected derivative residual $R$ gives
\[
 f(a)-\min_{[0,c]}f\le R^2/(2\mu).
\]
At zero, a nonnegative derivative has residual zero; at capacity, a nonpositive derivative has residual zero. Otherwise the absolute derivative bound is sufficient. This follows by minimizing the strong-convexity supporting quadratic over the feasible interval and, conservatively, over the line. Here $\mu=2$. Equation~\eqref{eq:cubic63}, its $41/32$ and $9/4$ coefficients, analytic shock moments, and outward derivatives are independently checked against rational formulas. Every method uses this kernel; it is not credited to the neural predictor.

\subsection{Work, prediction, and the economic object}
Let $A_r(\delta)$ be a bound on action endpoints at a query. Excluding training, a scalar operation account obeys
\[
 W_1(\delta)=O(dA_1),\qquad
 W_r(\delta)\le C d A_r(\delta)+
 A_r(\delta)q_r W_{r-1}(\delta/(4\beta)).
\]
The constants $dG_r,dM_r$ are bounded independently of $d$, so the declared endpoint and quadrature budgets need not grow exponentially with the state dimension. Primitive evaluations still cost $O(d)$. This removes the state lattice, not the horizon or shock-tree cost. The executed vectorization batches a query's innovation descendants. Its largest live state batch is reported; it does not claim the smaller memory of a fully streamed implementation. Increasing the horizon or the shock dimension may still be expensive. The new execution is scalar-control; the earlier full two-control result is retained rather than relabeled as a query-service experiment.

A useful prediction can reduce the endpoint count or return its own action from the initial partition. An inaccurate prediction can instead add work. This is a prediction-dependent algorithm, not a theorem of neural dominance. Learned warm starts have their own literature, including \citet{sakaue2023warm}; their discrete-convex setting differs from the original-law Bellman intervals here. Convex control policies and learned optimization parameters are studied by \citet{agrawal2020control}. Our comparison identifies the extra proposal cost and the resulting query transcript rather than claiming priority for learning a warm start.

The purchase is now a complete callable controller with a declared error contract. Its training cost, every subsequent certified query, and its return-time cost belong to the same service. A reduction in queried state storage is economically relevant for repeated decisions only after these costs are counted. The original cost normalization is maintained. Workload seconds and normalized expected-cost loss are reported separately; a price per unit of computation can convert the former to the latter only as an explicitly specified decision-maker parameter.

\section{Frozen neural continuations and policy-cost precision}\label{sec:frozen72}
The optimal-advantage estimator in Section~\ref{sec:centered69} makes a policy comparison informative, but its conditional optimal-value queries are expensive. A fitted continuation can replace those queries without changing the estimand. The necessary certificate concerns its Bellman residual across states, not its distance from an unknown optimizer. This places the original NBO continuation, rather than only a proposed action, inside the evaluation calculation.

\subsection{Residual oscillation is the relevant range}
Consider the bounded expected-cost model with common terminal cost $g$, discounts $\beta_t\ge0$, and a common initial-state law. Write $b_0=1$ and $b_t=\prod_{j<t}\beta_j$. Fix bounded measurable continuations $h_0,\ldots,h_T$, conditionally on their construction. Define
\begin{align*}
 Q_t^h(x,a)&=c_t(x,a)+\beta_tP_t^a h_{t+1}(x),\\
 r_t^h(x)&=\inf_{a\in A_t(x)}Q_t^h(x,a)-h_t(x),\\
 a_t^{h,\pi}(x)&=Q_t^h(x,\pi_t(x))-\inf_{a\in A_t(x)}Q_t^h(x,a).
\end{align*}
The infimum is over the original continuous feasible set. Suppose verified constants give $0\le a_t^{h,\pi}\le\eta_t^\pi$, and likewise for $\nu$, throughout the true state domain. Let $\omega_t$ bound the oscillation of $r_t^h$, and $\omega_T$ that of $g-h_T$. Set
\[
 \Omega_h=\sum_{t<T}b_t\omega_t+b_T\omega_T,
 \qquad H_h^\pi=\sum_{t<T}b_t\eta_t^\pi.
\]
These constants include acquisition, action rounding and arithmetic allowances wherever those operations affect the implemented policy or the verified functions. A small empirical residual alone is not such a certificate.

\begin{theorem}[Frozen-continuation comparison]\label{thm:frozen72}
Along the actual trajectory of $\pi$, put
\begin{equation}\label{eq:frozen-observation72}
 Z_h^\pi=\sum_{t<T}b_t\bigl[Q_t^h(X_t^\pi,\pi_t(X_t^\pi))-h_t(X_t^\pi)\bigr]
                 +b_T(g-h_T)(X_T^\pi).
\end{equation}
For any coupling of the two original trajectory laws, $D_h=Z_h^\pi-Z_h^\nu$ satisfies
\begin{align}
 \mathbb E D_h&=\mathbb E[J_0^\pi-J_0^\nu],\label{eq:frozen-identity72}\\
 -\Omega_h-H_h^\nu&\le D_h\le\Omega_h+H_h^\pi.\label{eq:frozen-range72}
\end{align}
Moreover, at every initial state,
\begin{equation}\label{eq:frozen-regret72}
 0\le J_0^\pi-V_0^*\le\Omega_h+H_h^\pi.
\end{equation}
Adding an arbitrary state-independent constant to each $h_t$ changes neither $D_h$ nor these oscillation bounds. With $h_t=V_t^*$, the result reduces to optimal-advantage centering.
\end{theorem}
\noindent The proof is in Supplement, Section~\ref{supp:frozen72}.

The terminal correction in \eqref{eq:frozen-observation72} is essential unless $h_T=g$. Conditional expectation, not pathwise equality with the realized cash-cost difference, supplies \eqref{eq:frozen-identity72}. The result uses the actual kernels and costs; it assumes neither a correct statistical transition model estimated from data nor a convex neural fit. The pointwise policy certificate also applies conditionally after an adaptive choice among valid continuations. The independent-sample confidence statement below requires its additional sampling conditions.

For $M$ conditionally fixed comparisons, suppose independent trajectory rows yield outward enclosures $[L_i,U_i]$ of $D_h$, intersected with its valid support. Write $R_h=2\Omega_h+H_h^\pi+H_h^\nu$. A simultaneous interval with failure probability at most $\alpha$ is
\begin{equation}\label{eq:frozen-ci72}
 \left[\overline L-R_h\sqrt{\frac{\log(2M/\alpha)}{2n}},\quad
       \overline U+R_h\sqrt{\frac{\log(2M/\alpha)}{2n}}\right].
\end{equation}
This follows by the bounded-sum inequality applied to both endpoint samples and a union bound. The numerical and sampling width allocations of Corollary~\ref{cor:precision69} apply with $R_h$. Fitting, uniform residual qualification, conditional integration and any subsequent selection are distinct paid operations. In particular, a cache that adapts across rows changes the controller's history: its rows cannot be treated as independent evaluations of one frozen Markov policy. A fixed cache or an independent episode-reset design is sufficient to retain the stated sampling construction.

\subsection{Conditional integration of an actual neural continuation}
A shallow rational affine--ReLU continuation has the form
\[
 h(y)=c+\ell^\top y+\sum_{j=1}^K v_j(w_j^\top y+u_j)_+.
\]
Its readout weights may have either sign. In the original investment laws the next state, conditional on $(x,a)$, is affine in one or two independent centered uniform innovations. A neuron therefore has projected center $z$ and nonnegative projection radii $A,B$. Write $z_+=\max\{0,z\}$.

\begin{proposition}[Exact conditional ReLU moments]\label{prop:moment72}
For independent uniform projections of radii $A,B>0$, the conditional means of a neuron are
\begin{align}
 \mathcal M_1(z;A)&=\frac{(z+A)_+^2-(z-A)_+^2}{4A},\label{eq:moment-one72}\\
 \mathcal M_2(z;A,B)&=\frac{1}{24AB}\bigl[(z+A+B)_+^3-(z+A-B)_+^3\nonumber\\
 &\hspace{37mm}-(z-A+B)_+^3+(z-A-B)_+^3\bigr].\label{eq:moment-two72}
\end{align}
A zero projection radius is removed; if all radii vanish the mean is $z_+$. Consequently $P^ah$ is a finite exact rational calculation for rational primitives, states, actions and weights. For $n$ trajectories and at most $K$ neurons per date, its arithmetic-operation count is $O(nTKd)$, exclusive of fitting, trajectory construction and uniform qualification. No optimal-value descendant query occurs in this calculation.
\end{proposition}
\noindent The proof is in Supplement, Section~\ref{supp:moments72}.

The operation count does not bound rational bit complexity. Near-zero projected radii can also make a naive floating evaluation unstable; directed rounding or rational evaluation, including the zero-radius branches, is required. The executed moment catalogue checks the rational identities against an independently written polygon integral. It is a conditional-integration experiment, not a new trained-policy cost comparison. The formulas do not imply that the achieved residual oscillation of an arbitrary fitted network is small.

\subsection{Qualifying a fitted continuation without a Bellman tree}
Uniform qualification need not query $V^*$. Suppose $Q_t^h$ is Lipschitz in state and action with constants $L_{x,t},L_{a,t}$, and $h_t$ has constant $L_{h,t}$, all in $\ell^1$. In the investment capacity simplex, an action feasible at $y$ can be repaired continuously at $x$ with displacement at most $\kappa\|x-y\|_1$, where $\kappa=1/(8d)$. Thus a valid residual Lipschitz constant is
\[
 L_{r,t}=L_{x,t}+\kappa L_{a,t}+L_{h,t}.
\]
Let a finite state cover have fill radius $\rho$, and let $[l_j,u_j]$ enclose $r_t^h$ at every cover point. Then
\begin{equation}\label{eq:residual-cover72}
 \min_jl_j-L_{r,t}\rho\le r_t^h(x)
       \le\max_ju_j+L_{r,t}\rho
\end{equation}
throughout the cube. At a cover point, a feasible action mesh of fill radius $s$ gives a lower infimum bound by subtracting $L_{a,t}s$ from the smallest action-value lower endpoint; the smallest upper endpoint is an upper bound. Exact moments can supply those action endpoints. None of these steps requires convexity of $h$ or a true-value curvature minorant for $Q^h$.

For $N_t$ state-cover points and $A_t$ actions per point, this qualification costs $O(\sum_tN_tA_tK_{t+1}d)$ arithmetic operations, plus certification of the moduli and any actor-specific greedy-gap bound. Its cover can still be expensive in high dimension. The result replaces a recursive optimal-value oracle by an explicit approximation-and-qualification tradeoff; it does not make that qualification free. The residual-cover proof, the action-mesh argument and the terminal allowance are given in the supplement. Value-function control variates have an established variance-reduction literature \citep{munos2006variance}; the additional requirement here is an original-law, whole-domain certificate for the fitted NBO continuation and the actual acquired policy.

\section{Fitting a continuation for qualification}\label{sec:qualification73}
A continuation used to evaluate a policy should be fitted to the error that determines the evaluation cost. The residual-cover bound in \eqref{eq:residual-cover72} supplies such an objective. We now make its optimization constructive. The future continuation is held fixed at each backward step; neither an optimal-value oracle nor success of unrestricted hidden-layer training is assumed.

Fix a date and suppress its subscript. Let $x_1,\ldots,x_N$ be a state cover of radius $\rho$ in the $\ell^1$ metric, and suppose
\[
 y_j^-\le (\mathcal T h_+)(x_j)\le y_j^+.
\]
These are verified bounds on the original continuous-action infimum, not only on a proposal menu. Write $h_\theta=\sum_{k=1}^K\theta_k\phi_k$, where the fixed features have verified Lipschitz constants $L_k$ and $|\theta_k|\le B_k$. The constant feature is allowed with $L_k=0$. Let $C$ be a verified Lipschitz constant of $\mathcal T h_+$, including state-dependent feasibility repair. With $F_{jk}=\phi_k(x_j)$, define
\begin{equation}\label{eq:readout-objective73}
 \mathcal W(\theta)=\max_j(y_j^+-F_j\theta)
       -\min_j(y_j^--F_j\theta)
       +2\rho\left(C+\sum_k L_k|\theta_k|\right).
\end{equation}

\begin{theorem}[Qualification-optimal readout]\label{thm:readout73}
For every feasible $\theta$, $\operatorname{osc}(\mathcal T h_+-h_\theta)\le\mathcal W(\theta)$. Minimizing \eqref{eq:readout-objective73} is a finite linear program with $2K+2$ variables and $2N+2K$ inequality constraints, apart from variable bounds. For probability vectors $p,q$ on the nodes, set $\lambda_k=2\rho L_k$ and
\begin{equation}\label{eq:readout-dual73}
 \mathcal D(p,q)=p^\top y^+-q^\top y^-+2\rho C
   -\sum_k B_k\bigl(|[F^\top(q-p)]_k|-\lambda_k\bigr)_+.
\end{equation}
Then $\mathcal D(p,q)\le\inf_\theta\mathcal W(\theta)$, with equality for some $p,q$. Thus any feasible readout and exact probability vectors certify an achieved optimization gap of at most $\mathcal W(\theta)-\mathcal D(p,q)$.
\end{theorem}

The proof is in the Technical Supplement. It constructs the primal constraints and dual support function explicitly. Rationalizing a solver's proposed readout and probability vectors gives a checkable certificate; a convergence flag is not the certificate. Label intervals, features, moduli, coefficient bounds and the cover remain part of the numerical contract.

Applying the theorem backwards gives a fitted continuation at every date. Its achieved ranges $\mathcal W_t$, together with the terminal oscillation and the actual implemented greedy-gap bounds, can be substituted directly into Theorem~\ref{thm:frozen72}. This connects the attained fitting objective to both original-optimum regret and the range of the original-law cost observation. Each readout minimizes its stated one-step qualification bound conditional on the fixed future. It does not minimize the joint multi-date range, the true policy loss, or a hidden-feature objective. The next backward step recomputes its labels and moduli from the continuation actually returned.

The label construction in \eqref{eq:residual-cover72} uses exact conditional moments for a shallow ReLU future and a verified continuous-action cover. Its work is charged in addition to the linear program, coefficient rationalization, modulus calculation and dual check. A conventional basis with the same number of coefficients receives the same labels and the same objective; it uses its own valid feature moduli. This is a comparison of qualified approximations, not an assertion that different bases span the same functions or that a neural encoding has an intrinsic advantage.

\section{Sequential qualification of policy costs}\label{sec:sequential73}
A fixed-sample interval does not by itself justify repeated inspection until a controller appears acceptable. Nor do cost observations from changing controllers identify the last controller's cost. We extend the frozen-continuation identity to predictable episodes and distinguish these targets. This uses established time-uniform concentration principles \citep{howard2020time}; confidence sequences have also been used for gated deployment in contextual bandits \citep{karampatziakis2021}. Here the observations are verified, continuous-law NBO cost contrasts, not importance-weighted bandit rewards.

Let $\mathcal F_{i-1}$ contain all information before episode $i$. Conditional on it, choose two complete controllers, a common initial law, a continuation $h_i$ fixed through the episode, and valid coefficients of Theorem~\ref{thm:frozen72}. The resulting cost contrast $D_i$ has conditional mean $\Delta_i$ equal to the original-law expected cost difference and predictable support $[a_i,b_i]$. Let $[L_i,U_i]\subseteq[a_i,b_i]$ enclose $D_i$ for the entire continuous innovation cell. Write $V_n=\sum_{i\le n}(b_i-a_i)^2$. For fixed positive numbers $\lambda_j,w_j$, $1\le j\le J$, with $\sum_jw_j\le1$, define
\begin{equation}\label{eq:anytime-boundary73}
 B_n=\min_{1\le j\le J}
       \left\{\frac{\log(2/(\alpha w_j))}{\lambda_j}
                    +\frac{\lambda_j V_n}{8}\right\}.
\end{equation}

\begin{theorem}[Predictable original-law cost certificate]\label{thm:anytime73}
Under the preceding conditional-law, support and enclosure conditions, with probability at least $1-\alpha$, simultaneously for every $n\ge1$,
\begin{equation}\label{eq:anytime-interval73}
 \frac{\sum_{i\le n}L_i-B_n}{n}
 \ \le\ \frac1n\sum_{i\le n}\Delta_i
 \ \le\ \frac{\sum_{i\le n}U_i+B_n}{n}.
\end{equation}
Each endpoint may also be intersected with the corresponding average predictable support endpoint. The assertion holds at every almost surely finite stopping time. If the complete controller pair and initial-law/reset contract stay fixed, then $\Delta_i=\Delta$, even when the qualified continuation changes predictably between episodes. Otherwise the target is the average conditional cost difference, not the last controller's cost.
\end{theorem}

The proof applies conditional Hoeffding bounds to each endpoint and Ville's inequality to a finite collection of nonnegative supermartingales. It requires no independent or identically distributed deployment contexts. It does require the stated conditional transition laws. Within-episode cache adaptation is allowed only when the controllers are treated as history-dependent algorithms and the greedy-gap qualification covers every memory state that can occur. Repeating the same program identifier while carrying a different cache state does not establish a fixed conditional target.

The implementation uses a finite dyadic collection of $\lambda_j$ and exact rational upper bounds for the logarithms. These give conservative, computable versions of \eqref{eq:anytime-boundary73}. It reserves the support and target identity before accepting an episode. A failed numerical evaluation consumes the full previously qualified support interval and its actual work; it is not deleted after inspecting its outcome. A failure of the support certificate itself stops qualification.

\begin{corollary}[Costed replacement after optional stopping]\label{cor:adoption73}
For a fixed target pair, let $[l_N,u_N]$ be the interval at a finite stopping time $N$. Let a declared deployment count be $m\ge1$, a one-time replacement and qualification charge be $F\ge0$, and the incremental per-deployment charge be $c$, all in the same units as the policy cost. On the simultaneous event of Theorem~\ref{thm:anytime73}, the net incremental cost lies in
\[
 [l_N+F/m+c,\ u_N+F/m+c].
\]
Consequently $u_N+F/m+c<0$ certifies strict net improvement. Inclusion of $[l_N,u_N]$ in a declared equivalence band certifies equivalence at that band, not strict superiority. New adaptively selected controller pairs may be qualified in separate epochs with conditional error allocations whose sum is at most the overall error budget.
\end{corollary}

A wall-clock reading is not a monetary charge. Measured setup, validation, rejected evaluations and successful retrievals are retained separately; prices are explicit decision inputs. The number of qualification episodes is not automatically the number of deployments over which a setup cost is amortized. These distinctions let the original investment decision use a valid stopping rule without silently changing either its estimand or its units.

\section{From a constructed controller to an adoption decision}\label{sec:displacement75}
The attained readout objective controls a uniform residual account. It does not, by itself, identify the expected cost of the controller constructed from that readout. We now qualify that cost directly. A finite acquired controller supplies information that a global payoff bound ignores: two controllers with nearby actions cannot separate arbitrarily under a common innovation. This information remains available when both controllers are discontinuous.

\subsection{A computable pathwise range}
Let $\pi,\nu$ be fixed complete controllers, constant on finite observation partitions $\mathcal C_t^\pi,\mathcal C_t^\nu$ of a bounded invariant state domain. Closures are used when computing distances. Their actions are feasible on the entire corresponding cell. Suppose, uniformly over innovations and feasible pairs,
\begin{align*}
 \|f_t(x,a,z)-f_t(y,b,z)\|&\le L_{f,t}\|x-y\|+B_{f,t}\|a-b\|,\\
 |c_t(x,a)-c_t(y,b)|&\le L_{c,t}\|x-y\|+B_{c,t}\|a-b\|,
 \qquad |g(x)-g(y)|\le L_g\|x-y\|.
\end{align*}
All constants concern the original primitives. Write $b_t=\prod_{s<t}\beta_s$. For $r\ge0$, the finite action-displacement modulus is
\begin{equation}\label{eq:cell-modulus75}
 \alpha_t(r)=\max_{\operatorname{dist}(\overline C,\overline D)\le r}
                 \|\pi_t(C)-\nu_t(D)\|,
 \quad (C,D)\in\mathcal C_t^\pi\times\mathcal C_t^\nu.
\end{equation}
It is nondecreasing. Its value at zero need not vanish, even for a discontinuous controller compared with itself: adjacent closed cells can carry different actions. Exact controller identity is checked separately, on a common refinement, rather than inferred from this deliberately conservative maximum.

\begin{theorem}[Finite-controller displacement certificate]\label{thm:displacement75}
Couple the original trajectory laws by a common initial state and common innovations. Set $\delta_0=0$ and, with $D_X$ the state-domain diameter,
\[
 a_t=\alpha_t(\delta_t),\qquad
 \delta_{t+1}=\min\{D_X,L_{f,t}\delta_t+B_{f,t}a_t\}.
\]
If $G^\pi,G^\nu$ denote the two realized discounted costs, then
\begin{equation}\label{eq:displacement-support75}
 |G^\pi-G^\nu|\le K_{\pi\nu}:=
 \sum_{t<T}b_t(L_{c,t}\delta_t+B_{c,t}a_t)+b_TL_g\delta_T.
\end{equation}
Any independently established cost-range bound may reduce $K_{\pi\nu}$. If the complete controllers agree on a common refinement, their coupled costs agree exactly and the bound can be set to zero. No regularity of the acquired actors or access to a policy-value function is required. Direct enumeration uses $O(\sum_t|\mathcal C_t^\pi||\mathcal C_t^\nu|d)$ arithmetic work for rectangular cells in the $\ell^1$ metric; the resulting bound is qualified before cost sampling.
\end{theorem}

\begin{proof}
Suppose the state displacement at date $t$ is at most $\delta_t$. The actual cell pair has closure distance no greater than that displacement, so its action difference is at most $a_t$. The transition inequality and invariance give the next displacement bound. Induction begins at the common initial state. Applying the cost inequalities at every date and summing proves \eqref{eq:displacement-support75}. Identical controllers choose identical actions at identical states, and induction then gives identical trajectories. Rectangle distances are sums of coordinatewise interval distances; enumerating all pairs gives the stated work.
\end{proof}

For the original two-state investment economy, valid constants are
\[
 L_f=\frac{11}{16},\quad B_f=\frac34,\quad
 L_c=4,\quad B_c=\frac{13}{16},\quad L_g=5,\quad D_X=2.
\]
The supplement derives them directly from the bilinear transition, quartic adjustment cost and shortage term. The common innovation cancels only in the displacement inequality; both marginal trajectories retain the original continuous innovation law. The certificate is thus not a replacement of the stochastic economy by its conditional mean. Its finite pair computation is paid, and can remain expensive when observation partitions are large.

\subsection{Numerical width and economic selection}
Condition on the frozen controllers. For each of $M$ declared estimands $\mu=\mathbb E D$, let independent continuous-bin rows produce endpoints $L_i\le D_i\le U_i$ in a prequalified support $[a,b]$. Define $R=b-a$, the unbiased sample variances $s_L^2,s_U^2$, and $\ell\ge\log(4M/\gamma)$. The endpoint empirical-Bernstein interval is
\begin{equation}\label{eq:cost-interval75}
 [l,u]=[\overline L-r_L,\overline U+r_U]\cap[a,b],\qquad
 r_V=\sqrt{2s_V^2\ell/n}+\frac{7R\ell}{3(n-1)}.
\end{equation}
The standard bounded-sample inequality of \citet{maurer2009}, applied to both endpoints and all estimands, gives simultaneous coverage at least $1-\gamma$. Independence is between rows; different controllers may share every row. Before support clipping, the interval width is exactly
\begin{equation}\label{eq:width-ledger75}
 (\overline U-\overline L)
 +\sqrt{2s_L^2\ell/n}+\sqrt{2s_U^2\ell/n}
 +\frac{14R\ell}{3(n-1)}.
\end{equation}
The first term contains bin, acquisition-branching and arithmetic enclosures. Its separate realized contributions can be audited without claiming that a larger sample removes acquisition discontinuities. The exact integer calculation below records both sampling radii, the range term and mean bin width. Support clipping cannot enlarge the interval. A finite-sample interval is not made valid at arbitrary stopping times by inspecting it repeatedly; the predictable construction in Section~\ref{sec:sequential73} has its own assumptions.

For adoption, write $[l_k,u_k]$ for the gain $\mathbb E(G^0-G^k)$. Let $D$ be the intended number of future deployments, $\lambda\ge0$ the value of a unit of normalized loss, $K_k\ge0$ the priced one-off construction, qualification and installation cost, and $d_k$ the priced incremental cost per deployment. On the same coverage event the total net adoption benefit belongs to
\begin{equation}\label{eq:adoption75}
 [D(\lambda l_k-d_k)-K_k,\ D(\lambda u_k-d_k)-K_k].
\end{equation}
Consequently $D>K_k/(\lambda l_k-d_k)$ supports strictly beneficial adoption whenever the denominator is positive. A negative upper endpoint rules it out. This statement is simultaneous in the nonnegative economic prices and deployment counts; choosing them after inspection does not create an uncovered cost estimand. They remain external economic inputs, and the inference sample size $n$ is not the deployment count $D$. The preceding readout, feasibility, support and inference results now terminate in a decision concerning the actual returned NBO controller.

\section{Transporting purchased Bellman witnesses}\label{sec:reuse72}
A second way to control the descendant cost is to reuse a calculation already purchased. Exact memoization only recognizes the same state. The following rule transfers a lower Bellman bound and its feasible upper witness to a nearby state, after repairing capacity and accounting for the full displacement. It is compatible with either a neural or a conventional generator of the original certified record.

\subsection{An original-optimum transport certificate}
At remaining horizon $r$, a record at $y$ contains an action $a_y\in A(y)$ and numbers satisfying
\[
 L_y\le V_r^*(y)\le Q_r^*(y,a_y)\le U_y.
\]
It also identifies the economic primitives, remaining horizon and numerical contract under which the record was obtained. A file hash detects changed bytes; validity of the inequalities is a separate premise. For the original scalar and two-control economies, let
\[
 L_V=G_r,\qquad L_Q=\frac8d+\frac{11\beta}{16}G_{r-1},
 \qquad L_A=\frac{13}{16}+\frac{3\beta d}{8}G_{r-1}.
\]
These bound state variation of $V_r^*$ and state and action variation of $Q_r^*$. The action comparison may be made on the ambient simplex $a\ge0$, $\sum_i a_i\le1/4$: the original transition remains in the state cube there, even if a particular action violates the smaller capacity at one endpoint.

At a new state $x$, repair $a_y$ to an actual lattice action $\widetilde a\in A(x)$. Set $D=\|x-y\|_1$ and $\Delta_a=\|\widetilde a-a_y\|_1$. Define
\begin{equation}\label{eq:witness-transport72}
 L(x)=\max\{0,L_y-L_VD\},\qquad
 U(x)=U_y+L_QD+L_A\Delta_a.
\end{equation}
The truncation at zero uses the nonnegative costs of this economy. It is not part of a general signed-cost transfer rule.

\begin{theorem}[Reuse after capacity repair]\label{thm:reuse72}
Under the stated record and modulus conditions,
\[
 L(x)\le V_r^*(x)\le Q_r^*(x,\widetilde a)\le U(x).
\]
Hence outward endpoints with $\overline U(x)-\underline L(x)\le\delta$ permit the actual repaired action to be returned without a new Bellman query. An identity mismatch, failed feasibility check or excessive width requires a new verified calculation; proximity alone licenses no reuse.

The same assertion holds after any history-dependent selection of valid records. At each future date, suppose the selected record or fresh query supplies the original local-gap and acquisition premises of Theorem~\ref{thm:query63}. The resulting, possibly history-dependent controller retains that theorem's original-optimum loss sum. It need not coincide with the action returned by a fresh search at the same tolerance.
\end{theorem}
\noindent The proof is in Supplement, Section~\ref{supp:transport72}.

The implemented repair first rounds the new capacity downward on the $2^{-32}$ lattice, scales a vector action proportionally when necessary, and rounds each coordinate downward. The certificate charges its exact resulting $\ell^1$ displacement, not merely the change in capacity. Retrieval may use approximate distances or a learned index as an ordering heuristic. Every returned hit must still satisfy \eqref{eq:witness-transport72} under outward arithmetic.

\subsection{A finite-cover account and its premises}
A small number of hits does not establish a uniform reuse complexity bound. Such a bound requires a fill precision, a coverage argument and a retention rule. Let the action quantum be $q$, let $m\in\{1,2\}$, and suppose every fresh record is purchased with gap at most $\gamma\delta$, where $0<\gamma<1$. The repair just described obeys
\[
 \Delta_a\le\kappa D+(m+1)q,\qquad \kappa=1/(8d).
\]
Consider the exact-endpoint transport rule, or an outward rule with a known total additional width allowance $\xi\ge0$. Define
\begin{equation}\label{eq:reuse-radius72}
 \rho_*=
 \frac{(1-\gamma)\delta-L_A(m+1)q-\xi}
      {L_V+L_Q+\kappa L_A}.
\end{equation}

\begin{proposition}[Paid fills and complete reuse work]\label{prop:cover72}
Suppose $\rho_*>0$, all valid records are retained, the context is fixed, and every stored candidate is checked before a miss is declared. A record within distance $\rho_*$ certifies a hit. On the unit state cube, the number $M_r$ of fresh records at remaining horizon $r$ is at most
\begin{equation}\label{eq:packing72}
 M_r\le\left(1+\frac d{\rho_*}\right)^d.
\end{equation}
If $C_r(\gamma\delta)$ bounds the complete cost of a fresh fill, a workload of $B$ requests obeys
\begin{equation}\label{eq:reuse-work72}
 W_B\le W_{\mathrm{setup}}+M_rC_r(\gamma\delta)
       +\sum_{b=1}^B\bigl(W_b^{\mathrm{lookup}}+W_b^{\mathrm{check}}
                     +W_b^{\mathrm{load/store}}\bigr).
\end{equation}
The lookup and check costs include rejected candidates. Naive retrieval uses $O(M_rd)$ arithmetic operations per request, apart from rational bit costs. Storage, eviction, invalidation and context changes require their own additional fill accounts when the stated retention premises do not hold.
\end{proposition}
\noindent The proof is in Supplement, Section~\ref{supp:cover72}.

This bound makes the dimensional price of uniform reuse explicit. It is not a dimension-free caching theorem. It also separates a first expensive Bellman fill from later transport checks. Caching at the workload root can avoid all descendants on a hit, but does not accelerate the descendants of a miss. The executed bounded cache uses 128 records and the requested fill tolerance, not the tighter $\gamma\delta$ premise. Its operation counts below are direct observations, not an application of \eqref{eq:packing72}. A cache at recursive descendants would need a separate execution and ledger.

\subsection{Constructed continuations and attained qualification}\label{sec:study73}
The readout theorem is executed in four predeclared backward constructions in the original two-state, two-date, scalar-control economy. Two midpoint covers, with eight or sixteen cells per coordinate, are crossed with additive shallow ReLU and degree-four Bernstein features. Each basis has nine coefficients bounded by eight. At each date the future is the continuation actually constructed at the next date. Exact conditional moments supply thirty-three continuous feasible action values per node and thirty-three robust, downward-quantized acquired-action proposals. The original state-dependent capacity and continuous innovation law are unchanged.

\begin{table}[htbp]
\centering
\caption{Attained continuation qualification in the original investment economy}\label{tab:qualification73}
{\small\setlength{\tabcolsep}{4pt}
\begin{tabular}{rlrrr}
\hline
$n$ & Basis & $\Omega_h$ upper & Regret upper & Max. dual gap \\
\hline
8 & ReLU & 5.39764 & 8.46716 & 0.000000014 \\
8 & Bernstein & 5.36565 & 8.11941 & 0.000000000 \\
16 & ReLU & 3.73299 & 5.28939 & 0.000000014 \\
16 & Bernstein & 4.43474 & 5.81890 & 0.000000000 \\
\hline
\end{tabular}}
\par\smallskip
{\footnotesize Each row is one predeclared, complete two-date construction. Bounds are rounded upward; the exact fractions and dates are retained. Regret upper denotes the proved original-optimum account, not an estimated loss or an attained economic equivalence margin. The maximum dual gap concerns only the declared fixed-feature, fixed-future readout objectives.}
\end{table}


The table reports the attained two-date residual account $\Omega_h$, the original-optimum bound $\Omega_h+H_h^\pi$, and the largest exact one-step optimization gap. A smaller qualification objective is not itself an identified expected-cost gain. The matched least-squares ablation in the appendix uses each service's identical features and fixed-future labels; it tests the fitting objective, not a separately propagated controller. Both fitting and uniform qualification are paid, and the explicit state cover is tensor based.

All eight datewise optimization gaps are below $1.4\times10^{-8}$, against the predeclared tolerance $1/65536$. Their qualified residual bounds are strictly below the matched label-loss bounds; the reductions range from about 18 to 63 percent. The resulting global regret bounds remain conservative, as the table makes explicit. Those objective reductions do not establish a narrow economic equivalence interval or cheaper qualification at the earlier $1/256$ margin.

Every new action value is independently reconstructed using the retained original rational terminal formula or a different exact integration method for the fitted continuation. Continuous-action lower covers, whole-cell choices and primal--dual bounds are then recomputed. These constructions add actual qualified continuations, not new independent cost observations. The stopping theorem has separate exact finite-model regressions; the earlier fixed-sample equivalence intervals are not retroactively given an anytime interpretation.

\section{Direct qualification of the constructed NBO controllers}\label{sec:study75}
The four backward constructions in Section~\ref{sec:study73} supply complete, memoryless acquired controllers. Their continuation coefficients were trained for the verified oscillation objective, not for the cost sample used here. We now evaluate those exact returned controllers under the original continuous law. This closes a different part of the argument from the earlier query-proposal experiment: the fitted continuation determines the action tables, rather than merely suggesting optional probes to a common recursive search.

\subsection{A prospective, matched comparison}
The design, source files and four controller arrays were frozen before this experiment. It compares ReLU and Bernstein continuations with nine coefficients and the same coefficient box, original economic information, cover resolution and backward construction rules. Each basis is evaluated at $n=8$ and $n=16$ cells per coordinate, with two dates and action quantum $1/4096$. Both dates of each attained controller are retained. The feasible zero-investment controller supplies a common economically interpretable incumbent. It is not a substitute for the conventional baselines in the earlier studies.

The original initial state is uniform on $[0,1]^2$. At each date a fresh uniform innovation on $[-1/32,1/32]$ enters the two coordinates with opposite signs. The discount factor is $15/16$. Each of 131,072 predetermined rows represents its entire continuous product bin, with forty bits for each initial coordinate and each innovation. Every controller uses the same row. Thus there are 655,360 controller-path evaluations but only 131,072 joint cost rows. A fixed stream makes these rows reproducible; the independent-bin model is the statistical assumption.

There are thirteen declared estimands: five absolute costs, four constructed-controller-minus-zero differences, two ReLU-minus-Bernstein differences and two sixteen-cell-minus-eight-cell differences. Their simultaneous failure allowance is $1/100$. The paired width target and equivalence half-margin were separately set to $1/256$ before execution. All endpoints, all contrasts and the unsuccessful precision outcome are retained. Neither the sample count nor the stream was enlarged after observing a result.

\input{tables/cost75}
\input{tables/contrasts75}

Both ReLU controllers have lower expected cost than their resolution-matched Bernstein controllers on the simultaneous event. At sixteen cells, the ReLU-minus-Bernstein interval is $[-0.077238,-0.074459]$, rounded outward. Relative to the zero-investment incumbent, the same ReLU controller has gain at least $0.164484$ in normalized cost units. These are signed statements about actual complete-policy costs, not comparisons of training losses or uniform regret bounds.

The qualification criterion and the economic ordering need not coincide. At eight cells, the retained Bernstein construction has a smaller global oscillation certificate, yet the attained ReLU controller has lower original-law expected cost. The two quantities concern different risks: a uniform bound against the optimum and an initial-law average of the returned policy. The present evidence identifies the latter comparison, without retroactively reversing the earlier query-service timings or ranking the new table controllers against an unevaluated recursive FVI controller.

Refinement lowers cost within both bases. Those signed intervals also lie inside the declared equivalence band. Strictly lower cost and equivalence at a positive tolerance are compatible: the former concerns zero, and the latter concerns whether the entire difference is economically small at the declared numerical scale. The margin here is a predeclared numerical tolerance, not a calibrated monetary indifference threshold.

\subsection{Precision, independent recomputation and work}
Seven of the eight paired intervals meet the width target. The sixteen-cell ReLU-versus-zero interval has width approximately $0.00400629$, which exceeds $1/256$. Its gain nevertheless remains strictly positive. We report precision failure and signed identification separately, rather than widening the target to make every row pass. The absolute-cost intervals have the same simultaneous coverage but were not the objects of the paired-width target.

\input{tables/width75}

The range calculation precedes sampling and enumerates 206,080 cell pairs. An independent rational implementation reconstructs its datewise maxima. Production then encloses every continuous path bin using outward binary64 arithmetic. A separate implementation recomputes all 655,360 paths with arbitrary-precision dyadic integer arithmetic and an expanded, rather than factorized, transition. Final inference is calculated from the independently reconstructed endpoints using exact integer sums, rational variances and outward square-root bounds. The two implementations have no nonoverlapping path enclosures. Exact point and face fixtures supplement, rather than replace, the whole-bin calculations.

No sampled path bin crosses an acquired-action boundary in either implementation. The largest independent final path width is below $1.23\times10^{-11}$. Thus action ambiguity is zero in this realized catalogue, while the confidence width is dominated by sampling and the deterministic range allowance. This observation does not permit a global smoothness assumption on the actors; boundary fixtures and complete cell-pair qualification continue to include both sides of a boundary.

\input{tables/work75}

The resource account charges controller loading, finite support qualification, production path evaluation, complete independent recomputation, exact inference, durable outputs and process startup and shutdown. Earlier fitting and construction clocks remain separate measured observations. Combining them with a current qualification clock is not a jointly executed cold-start timing. Repeating the same random indices on a second machine is an arithmetic audit, not another independent statistical sample. The earlier adverse neural cold-start and reuse comparisons remain unchanged in the Numerical and Development Appendix.

\subsection{The adoption region}
For the sixteen-cell ReLU controller, equation~\eqref{eq:adoption75} supports adoption over the zero-investment incumbent whenever
\[
 D(0.164484\lambda-d_k)>K_k.
\]
The analogous conservative threshold for replacing the resolution-matched Bernstein controller uses $0.074459$ instead. These lower bounds are rounded downward. They define a nonempty region of qualification prices, installation charges and deployment counts for which the adopted controller has positive covered net benefit. The full intervals also determine the region of ruled-out purchase and the intervening unresolved region.

This is an economic decision in the original investment model with an identified policy-cost term. It is not a calibrated welfare estimate or a claim that the neural method is universally cheaper to construct or deploy. Its contribution is the completed chain: attained continuation fitting, feasible implemented actions, a primitive-derived concentration range, an independently checked original-law cost comparison and an explicit priced adoption condition. The other original NBO applications remain available under their separate assumptions; the new empirical ordering concerns the four controllers and the law specified here.

\section{Conclusion}\label{sec:conclusion65}
Neural Bellman Operators connect the construction of a continuation to feasible decisions under the original economic law. The witness-producing implementation has an explicit own-future error account. Qualification-optimal readouts give an attained fitting objective and an exact optimization certificate; original-law moments and uniform residual bounds connect that continuation to the implemented policy. The finite-controller displacement theorem additionally turns the complete acquired action tables into a prequalified path-cost range without imposing actor smoothness.

The resulting direct comparison identifies lower expected cost for both attained ReLU controllers than for their matched Bernstein controllers in the original two-date investment economy. It retains one failed width target, all comparative intervals and the full cost of independent qualification. The signed gains support a nonempty priced adoption region. This is a claim about the actual frozen controllers and law, not a universal neural ranking, a calibrated welfare estimate or an inferred reversal of the earlier complete-work evidence.

The other resource distinctions remain important. Exact conditional moments avoid an optimal-value descendant tree when evaluating a qualified shallow continuation, while fitting and uniform qualification remain paid. Transported witnesses avoid descendants on a certified hit, while the first fill and every miss retain their horizon cost. Predictable cost intervals permit valid stopping only for their declared conditional target. The original NBO objective is advanced by a complete attained-fit--feasibility--original-law-cost--adoption chain, with prior applications and unfavorable comparisons preserved for continued assessment.

\bibliographystyle{ecta-fullname}
\bibliography{references}
\end{document}
